% ================= Topic 3 (EN) -- part b: problems 16--30 =================
\providecommand{\cpTooth}[3]{% #1 rows, #2 columns, #3 unit (cm): grid of toothpicks
  \foreach \j in {0,...,#1}{\foreach \i in {1,...,#2}{
    \draw[tmOrange!85!black,line width=1.6pt,line cap=round]
      ({(\i-1)*#3+0.05},{\j*#3})--({\i*#3-0.05},{\j*#3});}}
  \foreach \j in {1,...,#1}{\foreach \i in {0,...,#2}{
    \draw[tmOrange!85!black,line width=1.6pt,line cap=round]
      ({\i*#3},{(\j-1)*#3+0.05})--({\i*#3},{\j*#3-0.05});}}}

\begin{probbox}{16}{\relax}
On the dot grid shown, Jane stands at dot $A$. She tosses a fair coin to decide how to move: if the coin shows
\emph{heads} she moves \emph{up} one dot; if it shows \emph{tails} she moves \emph{right} one dot. After exactly four
tosses, Jane will stop at one of the five dots $P$, $Q$, $R$, $S$, $T$. What is the probability that Jane stops at dot $R$?
\begin{figC}\begin{scope}[scale=0.8]
  \foreach \y in {0,...,4}{\pgfmathtruncatemacro\m{4-\y}
    \foreach \x in {0,...,\m}{\fill[Navy] (\x,\y) circle (2.6pt);}}
  \node[below left] at (0,0) {$A$};
  \node[above right] at (0,4) {$P$};\node[above right] at (1,3) {$Q$};
  \node[above right] at (2,2) {$R$};\node[above right] at (3,1) {$S$};
  \node[above right] at (4,0) {$T$};
  \fill[tmOrange] (0,0) circle (3.4pt);
\end{scope}\end{figC}
\choices{$\dfrac12$}{$\dfrac38$}{$\dfrac9{16}$}{$\dfrac7{16}$}{$\dfrac5{16}$}
\end{probbox}
\probsol{
\textbf{Understanding the problem.} At each toss Jane moves \emph{up} ($\uparrow$, for heads) or \emph{right}
($\rightarrow$, for tails). After $4$ tosses she has made exactly $4$ moves. We need the probability that after these
$4$ moves she is at $R$.

\step{1}{Count all possible outcomes}
Each toss has $2$ outcomes (heads or tails), so four tosses in a row give
$2\times2\times2\times2=16$ sequences, for example $\uparrow\rightarrow\rightarrow\uparrow$.
Since the coin is fair, these $16$ sequences are \emph{equally likely}.
\begin{recallbox}\emph{Multiplication principle:} if the first task can be done in $a$ ways, the second in $b$ ways, \dots{} then
doing them one after another can be done in $a\times b\times\cdots$ ways. The \emph{probability} of an event (when the outcomes
are equally likely) $=\dfrac{\text{number of favourable outcomes}}{\text{number of possible outcomes}}$.\end{recallbox}

\step{2}{The final dot depends only on the number of up-moves}
If $k$ of the $4$ moves are up, then $4-k$ moves are to the right. The order of the moves does not change the final dot.
Starting from $A$: $4$ ups $\to P$; $3$ ups $\to Q$; $2$ ups $\to R$; $1$ up $\to S$; $0$ ups $\to T$.
So Jane stops at $R$ exactly when she makes \textbf{exactly $2$ moves up and $2$ moves right}.

\step{3}{List the favourable sequences}
We choose which two of the four turns (turn 1, 2, 3, 4) are $\uparrow$:
\[\uparrow\uparrow\rightarrow\rightarrow,\quad \uparrow\rightarrow\uparrow\rightarrow,\quad
\uparrow\rightarrow\rightarrow\uparrow,\quad \rightarrow\uparrow\uparrow\rightarrow,\quad
\rightarrow\uparrow\rightarrow\uparrow,\quad \rightarrow\rightarrow\uparrow\uparrow.\]
There are $6$ sequences (choosing $2$ turns out of $4$: $\dfrac{4\times3}{2}=6$ ways).

\step{4}{Compute the probability}
\[P(\text{stop at } R)=\frac{6}{16}=\frac38.\]

\step{5}{A picture: write the number of paths at each dot}
At each dot write the number of move sequences leading from $A$ to that dot. A dot can only be reached from the dot just to its
left or the dot just below it, so the number at each dot $=$ the sum of the numbers to its left and below (just like Pascal's triangle):
\begin{figC}\begin{scope}[scale=0.75]
  \foreach \x/\y/\n in {0/0/1,1/0/1,2/0/1,3/0/1,4/0/1,0/1/1,1/1/2,2/1/3,3/1/4,0/2/1,1/2/3,2/2/6,
     0/3/1,1/3/4,0/4/1}{
    \fill[tmBlue!12] (\x,\y) circle (9pt);
    \node[font=\small\bfseries,text=Navy] at (\x,\y) {\n};}
  \draw[AnsRed,thick] (2,2) circle (10pt);
  \node[right=10pt,text=AnsRed,font=\small] at (2,2) {$R$};
  \node[right=10pt,font=\small] at (4,0) {$T$};\node[right=10pt,font=\small] at (3,1) {$S$};
  \node[right=10pt,font=\small] at (1,3) {$Q$};\node[right=10pt,font=\small] at (0,4) {$P$};
\end{scope}\end{figC}
Along the last diagonal: $1+4+6+4+1=16$ (exactly the total number of sequences), and the number at $R$ is $6$.
\begin{warnbox}The five dots $P,Q,R,S,T$ are \emph{not} equally likely! Do not answer ``one of $5$ dots, so the probability is
$\frac15$''. The middle dot ($R$) can be reached in the most ways; each of the two end dots ($P$ and $T$) can be reached in only one way.\end{warnbox}
\verify{The probabilities of stopping at $P,Q,R,S,T$ are $\frac1{16},\frac4{16},\frac6{16},\frac4{16},\frac1{16}$;
they add up to $\frac{16}{16}=1$, as they should.}
The value $\frac38=\frac6{16}$ is option \textbf{(B)}.
\ans{(B) $\dfrac38$}
}

\begin{probbox}{17}{\relax}
Identical toothpicks are used to build rectangular grids of squares, as shown below. For example, the $1\times10$ grid
($1$ row of $10$ squares) uses $31$ toothpicks in total. How many toothpicks are needed for a $43\times10$ grid
($43$ rows, each with $10$ squares)?
\begin{figC}
  \begin{scope}[shift={(0,2.4)}]\cpTooth{1}{10}{0.42}\node[right] at (4.4,0.21) {$1\times10$};\end{scope}
  \begin{scope}[shift={(0,1.05)}]\cpTooth{2}{10}{0.42}\node[right] at (4.4,0.42) {$2\times10$};\end{scope}
  \begin{scope}[shift={(0,-0.6)}]\cpTooth{3}{10}{0.42}\node[right] at (4.4,0.63) {$3\times10$};\end{scope}
\end{figC}
\choices{$913$}{$860$}{$871$}{$903$}{$946$}
\end{probbox}
\probsol{
\textbf{Understanding the problem.} The $43\times10$ grid has $43$ rows of squares, $10$ squares in each row. The toothpicks lie in
two directions: \emph{horizontal} and \emph{vertical}. We count each kind separately and add (addition principle).

\step{1}{Check the method on the $1\times10$ grid}
\emph{Horizontal toothpicks:} the grid has $2$ horizontal lines (top edge and bottom edge), each made of $10$ toothpicks:
$2\times10=20$ toothpicks.
\emph{Vertical toothpicks:} there are $11$ vertical lines ($10$ squares need $11$ walls), each made of $1$ toothpick:
$11\times1=11$ toothpicks. Total: $20+11=31$, exactly as stated.
\begin{recallbox}A row of $10$ squares needs $10+1=11$ vertical walls (like the fence-post problem: $10$ gaps need $11$ posts).
In the same way, $43$ rows of squares have $43+1=44$ horizontal lines.\end{recallbox}

\step{2}{Count the horizontal toothpicks of the $43\times10$ grid}
There are $43+1=44$ horizontal lines, each $10$ toothpicks long: $44\times10=440$ toothpicks.

\step{3}{Count the vertical toothpicks of the $43\times10$ grid}
There are $10+1=11$ vertical lines, each $43$ toothpicks long (one toothpick per row): $11\times43=473$ toothpicks.

\step{4}{Add}
$440+473=913$ toothpicks.

\step{5}{Second method: add one row at a time}
The $1\times10$ grid has $31$ toothpicks. Each time a new row is placed on top, we need $10$ more horizontal toothpicks
(the top edge of the new row) and $11$ vertical ones, that is $21$ toothpicks (the bottom edge is already there). Going from
$1$ row to $43$ rows means adding $42$ rows:
\[31+42\times21=31+882=913.\]
\begin{warnbox}Do not compute $43\times31=1333$: when rows are stacked, two neighbouring rows \emph{share} a line of $10$ horizontal
toothpicks, so multiplying like that counts those toothpicks twice.\end{warnbox}
\verify{By the formula $(\text{rows}+1)\times10+11\times\text{rows}$: the $2\times10$ grid needs
$3\times10+11\times2=52$ toothpicks and the $3\times10$ grid needs $4\times10+11\times3=73$; counting directly in the picture also gives
$52$ and $73$. Both methods give the same result, $913$.}
So the answer is \textbf{(A)}.
\ans{(A) $913$ toothpicks}
}

\begin{probbox}{18}{\relax}
In the cafeteria at Braydon's school, the meats are beef and chicken, and the fruits are apple, pear and banana. Braydon is given a
random lunch consisting of one meat and one fruit. The probability that Braydon's lunch includes a banana is
\choices{$\dfrac13$}{$\dfrac23$}{$\dfrac12$}{$\dfrac15$}{$\dfrac35$}
\end{probbox}
\probsol{
\textbf{Understanding the problem.} A lunch $=$ (one meat, one fruit). The lunch is chosen at random, so all possible lunches are
equally likely. We need the probability that the lunch includes a banana.

\step{1}{List all possible lunches}
By the multiplication principle there are $2\times3=6$ lunches. We make a table (rows: meat, columns: fruit):
\begin{center}\begin{tabular}{|c|c|c|c|}\hline
 & Apple & Pear & \cellcolor{tmOrange!20}Banana\\\hline
Beef & Beef--Apple & Beef--Pear & \cellcolor{tmOrange!20}Beef--Banana\\\hline
Chicken & Chicken--Apple & Chicken--Pear & \cellcolor{tmOrange!20}Chicken--Banana\\\hline
\end{tabular}\end{center}

\step{2}{Count the lunches with a banana}
The ``Banana'' column has $2$ lunches: Beef--Banana and Chicken--Banana.

\step{3}{Compute the probability}
\[P(\text{banana})=\frac{2}{6}=\frac13.\]
\begin{recallbox}To simplify a fraction, divide the numerator and the denominator by their greatest common divisor:
$\dfrac26=\dfrac{2\div2}{6\div2}=\dfrac13$.\end{recallbox}
\begin{tipbox}The meat has no effect on the fruit. A quicker way to think: the fruit is one of $3$ equally likely kinds, and banana is
$1$ of the $3$ kinds, so the probability is $\frac13$.\end{tipbox}
\begin{warnbox}Do not add ``$2$ meats $+\,3$ fruits $=5$'' and answer $\frac15$: the number of lunches is $2\times3=6$, not $2+3=5$.\end{warnbox}
\verify{The probabilities of getting an apple, a pear and a banana are all $\frac13$; they add up to $\frac13+\frac13+\frac13=1$.}
So the answer is \textbf{(A)}.
\ans{(A) $\dfrac13$}
}

\begin{probbox}{19}{\relax}
A special six-sided die is made, with one number written on each face (numbers may repeat). When this die is rolled, the probability
of getting a multiple of $3$ is $\frac12$, and the probability of getting an even number is $\frac13$.
The six numbers on the faces could be
\choicesII{$1,2,3,5,5,6$}{$1,2,3,3,5,6$}{$1,2,3,4,6,6$}{$1,2,3,3,4,6$}{$2,3,3,3,5,6$}
\end{probbox}
\probsol{
\textbf{Understanding the problem.} The die has $6$ equally likely faces. From the two given probabilities we work out the
\emph{number of faces} showing a multiple of $3$ and the number of faces showing an even number, then check each option.

\step{1}{Turn probabilities into numbers of faces}
\begin{recallbox}With $6$ equally likely faces: $P=\dfrac{\text{number of favourable faces}}{6}$, so
the number of favourable faces $=P\times6$.\end{recallbox}
Faces with a multiple of $3$: $\frac12\times6=3$ faces.\quad Even faces: $\frac13\times6=2$ faces.

So we need: \textbf{exactly $3$ faces that are multiples of $3$} and \textbf{exactly $2$ even faces}.

\step{2}{Check each option}
Among the numbers $1$ to $6$: the multiples of $3$ are $3$ and $6$; the even numbers are $2$, $4$, $6$.
Note that $6$ is both even and a multiple of $3$, so it must be counted in \emph{both} columns.
\begin{center}\begin{tabular}{|c|c|c|c|c|}\hline
Option & Faces & Multiples of $3$ & Even faces & OK?\\\hline
(A) & $1,2,3,5,5,6$ & $3,6$ \;($2$ faces) & $2,6$ \;($2$ faces) & No\\\hline
\rowcolor{tmGreen!12}(B) & $1,2,3,3,5,6$ & $3,3,6$ \;($3$ faces) & $2,6$ \;($2$ faces) & Yes\\\hline
(C) & $1,2,3,4,6,6$ & $3,6,6$ \;($3$ faces) & $2,4,6,6$ \;($4$ faces) & No\\\hline
(D) & $1,2,3,3,4,6$ & $3,3,6$ \;($3$ faces) & $2,4,6$ \;($3$ faces) & No\\\hline
(E) & $2,3,3,3,5,6$ & $3,3,3,6$ \;($4$ faces) & $2,6$ \;($2$ faces) & No\\\hline
\end{tabular}\end{center}
\begin{warnbox}If you forget that $6$ is even, you will count only $2$ even faces for (D) and wrongly choose (D).
A number can belong to several groups at once!\end{warnbox}
\verify{For (B): $P(\text{multiple of }3)=\frac36=\frac12$ and $P(\text{even})=\frac26=\frac13$; both conditions hold.}
Only option \textbf{(B)} works.
\ans{(B) $1,2,3,3,5,6$}
}

\begin{probbox}{20}{\stars{1}}
Three bus routes run along the roads shown in the figure, all passing through stop $X$ (segments with the same number of small marks
are equal: $PX=XS$, $QX=XT$, $RX=XU$).
\begin{itemize}[leftmargin=1.2em,itemsep=1pt,topsep=2pt]
\item Bus $A$ runs the route $P\to X\to S\to X\to P$; each round trip takes $12$ minutes.
\item Bus $B$ runs the route $Q\to X\to T\to X\to Q$; each round trip takes $20$ minutes.
\item Bus $C$ runs the route $R\to X\to U\to X\to R$; each round trip takes $28$ minutes.
\end{itemize}
At 1:00 p.m., buses $A$, $B$, $C$ leave $P$, $Q$, $R$ respectively. Each bus travels at a constant speed, turns around instantly
at both ends of its route and runs without stopping until 11:00 p.m. Between 5:00 p.m. and 10:00 p.m., at how many moments are two
or more buses at stop $X$ at the same time?
\begin{figC}\begin{scope}[scale=0.9]
  \draw[Ink,thick] (0,1.6)--(0,-1.6);\draw[Ink,thick] (-1.9,0)--(1.9,0);\draw[Ink,thick] (-2,-1.35)--(2,1.35);
  \foreach \p in {(0,1.6),(0,-1.6),(-1.9,0),(1.9,0),(-2,-1.35),(2,1.35),(0,0)}{\fill[Navy] \p circle (2pt);}
  \node[above] at (0,1.6) {$P$};\node[below] at (0,-1.6) {$S$};\node[left] at (-1.9,0) {$Q$};
  \node[right] at (1.9,0) {$T$};\node[below left] at (-2,-1.35) {$R$};\node[above right] at (2,1.35) {$U$};
  \node[below right] at (0.05,-0.02) {$X$};
  \draw[tmBlue,thick] (-0.12,0.8)--(0.12,0.8);\draw[tmBlue,thick] (-0.12,-0.8)--(0.12,-0.8);
  \foreach \s in {-1,1}{\draw[tmOrange,thick] (\s*0.95-0.04,-0.12)--(\s*0.95-0.04,0.12);
                        \draw[tmOrange,thick] (\s*0.95+0.04,-0.12)--(\s*0.95+0.04,0.12);}
  \foreach \s in {-1,1}{\foreach \d in {-0.09,0,0.09}{
     \draw[tmGreen!80!black,thick] ($(\s*1,\s*0.675)+(\d*0.829,\d*0.559)+(0.067,-0.1)$)--++(-0.134,0.2);}}
\end{scope}\end{figC}
\choices{$18$}{$19$}{$20$}{$21$}{$22$}
\end{probbox}
\probsol{
\textbf{Understanding the problem.} We measure time as \emph{the number of minutes after 1:00 p.m.} and call this number $t$.
We need the values of $t$ at which at least two buses are at $X$, with $t$ between 5:00 p.m. and 10:00 p.m.

\step{1}{Convert the time window into minutes}
5:00 p.m. is $4$ hours after 1:00 p.m.: $t=4\times60=240$. \quad 10:00 p.m. is $9$ hours after 1:00 p.m.: $t=9\times60=540$.
So we look at $240\le t\le540$.

\step{2}{When is bus $A$ at $X$?}
One round trip of bus $A$ consists of $4$ equal legs: $P\to X$, $X\to S$, $S\to X$, $X\to P$ (because $PX=XS$).
The bus moves at constant speed, so each leg takes $12\div4=3$ minutes. Bus $A$ reaches $X$ at the end of the first leg and at the end
of the third leg of every round trip, that is at
\[t=3,\ 9,\ 15,\ 21,\ 27,\ \dots\]
These are exactly the \textbf{odd multiples of $3$} ($3\times1$, $3\times3$, $3\times5$, \dots). At the end of the second and fourth legs
the bus is at $S$ or $P$, not at $X$.

\step{3}{The same for buses $B$ and $C$}
Bus $B$: each leg takes $20\div4=5$ minutes, so it is at $X$ when $t=5,15,25,35,\dots$ -- the \textbf{odd multiples of $5$}.\\
Bus $C$: each leg takes $28\div4=7$ minutes, so it is at $X$ when $t=7,21,35,49,\dots$ -- the \textbf{odd multiples of $7$}.

\step{4}{When do two buses meet at $X$?}
\begin{recallbox}A number divisible by both $3$ and $5$ is divisible by $\text{lcm}(3,5)=15$. Similarly:
divisible by both $3$ and $7$ $\Leftrightarrow$ divisible by $21$; divisible by both $5$ and $7$ $\Leftrightarrow$
divisible by $35$.\end{recallbox}
Buses $A$ and $B$ are both at $X$ when $t$ is an odd multiple of $3$ and also an odd multiple of $5$, that is when $t$ is an
\textbf{odd number divisible by $15$} (an odd multiple of $15$). Similarly: $A$ and $C$ meet when $t$ is an odd multiple of $21$;
$B$ and $C$ meet when $t$ is an odd multiple of $35$.
\begin{proofbox}An odd multiple of $15$ is an odd number divisible by $15$, so it is an odd number divisible by $3$ (an odd multiple of $3$)
and an odd number divisible by $5$ (an odd multiple of $5$). Conversely, a number that is an odd multiple of $3$ and an odd multiple of $5$ is
odd and divisible by $15$.\end{proofbox}

\step{5}{List them for $240\le t\le540$}
\begin{center}\begin{tabular}{|l|l|c|}\hline
Two buses & Values of $t$ (odd multiples) & Count\\\hline
$A$ and $B$ ($15$) & $255, 285, 315, 345, 375, 405, 435, 465, 495, 525$ & $10$\\\hline
$A$ and $C$ ($21$) & $273, 315, 357, 399, 441, 483, 525$ & $7$\\\hline
$B$ and $C$ ($35$) & $245, 315, 385, 455, 525$ & $5$\\\hline
\end{tabular}\end{center}
(For example, in the first row: $15\times17=255$ is the first odd multiple greater than $240$; after that we keep adding $30$, because
two consecutive odd multiples of $15$ are $2\times15=30$ apart.)

\step{6}{Remove the moments counted more than once}
Adding the three rows gives $10+7+5=22$, but the moments $315$ and $525$ appear in \emph{all three rows}
(they are odd multiples of $3\times5\times7=105$: at those moments all three buses are at $X$). Each of these moments is counted $3$ times
but should be counted only once, so we subtract $2+2=4$:
\[22-4=18 \text{ moments.}\]
\begin{warnbox}The question asks for the number of \emph{moments}, not the number of ``pairs of buses meeting''. When all three buses are at
$X$ together, that is just \emph{one} moment. Adding $10+7+5=22$ and choosing (E) is wrong.\end{warnbox}
\verify{In clock times: 5:05, 5:15, 5:33, 5:45, 6:15, 6:45, 6:57, 7:15, 7:25, 7:39, 7:45, 8:15, 8:21, 8:35, 8:45, 9:03, 9:15, 9:45 p.m.
-- exactly $18$ moments, all between 5:00 p.m. and 10:00 p.m.}
So the answer is \textbf{(A)}.
\ans{(A) $18$}
}

\begin{probbox}{21}{\relax}
Each face of a cube is painted with exactly one colour. What is the smallest number of colours needed to paint the cube so that any
two faces sharing an edge have different colours?
\choices{$2$}{$3$}{$4$}{$5$}{$6$}
\end{probbox}
\probsol{
\textbf{Understanding the problem.} Two faces ``sharing an edge'' are adjacent faces. Two opposite faces (such as the top and the
bottom) have no common edge, so they may have the same colour. We must (1) show a colouring that uses few colours, and
(2) prove that fewer colours are impossible.

\step{1}{Two colours are not enough}
\begin{proofbox}Look at any corner of the cube, for example the top--front--right corner. Three faces meet there: the top, the front and the
right face. Any two of these three faces share an edge, so the three faces must have three different colours. Hence at least $3$ colours
are needed.\end{proofbox}

\step{2}{Three colours are enough}
Paint \emph{each pair of opposite faces the same colour}: top and bottom blue, front and back orange, left and right green.
Each face is adjacent to $4$ faces, which are exactly the four faces \emph{not} opposite to it, and these four faces carry the other two
colours. So no two adjacent faces have the same colour.
\begin{figC}
  \begin{scope}
   \fill[tmOrange!55,draw=Ink,thick] (0,0)--(1.6,0)--(1.6,1.6)--(0,1.6)--cycle;
   \fill[tmBlue!55,draw=Ink,thick] (0,1.6)--(1.6,1.6)--(2.4,2.3)--(0.8,2.3)--cycle;
   \fill[tmGreen!55,draw=Ink,thick] (1.6,0)--(2.4,0.7)--(2.4,2.3)--(1.6,1.6)--cycle;
   \node at (0.8,0.8) {front};\node at (1.2,1.95) {top};\node at (2,1.15) {right};
  \end{scope}
  \begin{scope}[shift={(4,0.1)},scale=0.5]
   \foreach \x/\y/\c in {1/2/tmBlue,0/1/tmGreen,1/1/tmOrange,2/1/tmGreen,3/1/tmOrange,1/0/tmBlue}{
     \fill[\c!55,draw=Ink,thick] (\x,\y) rectangle ++(1,1);}
   \node[font=\scriptsize] at (2,-0.45) {net of the cube};
  \end{scope}
\end{figC}

\step{3}{Conclusion}
At least $3$ colours are needed (step 1) and $3$ colours are enough (step 2), so the smallest number of colours is $3$.
\begin{warnbox}Finding one colouring with $3$ colours is \emph{not enough} to claim that $3$ is the smallest; you must explain why $2$ colours
fail (the three faces at a corner). On the other hand, thinking ``$6$ faces need $6$ colours'' forgets that opposite faces may share a
colour.\end{warnbox}
\verify{On the net, each square shares edges only with squares of other colours (when folded, two squares of the same colour become
opposite faces).}
So the answer is \textbf{(B)}.
\ans{(B) $3$ colours}
}

\begin{probbox}{22}{\relax}
How many different pairs $(m,n)$ can be formed from the integers $1,2,3,\dots,20$ such that $m<n$ and $m+n$ is even?
\choices{$55$}{$90$}{$140$}{$110$}{$50$}
\end{probbox}
\probsol{
\textbf{Understanding the problem.} We choose two different numbers out of $20$; because $m<n$, each pair of two numbers is counted
\emph{only once} (the smaller one is $m$, the larger one is $n$). Condition: the sum $m+n$ is even.

\step{1}{When is the sum of two numbers even?}
\begin{recallbox}even $+$ even $=$ even;\quad odd $+$ odd $=$ even;\quad even $+$ odd $=$ odd.\end{recallbox}
So $m+n$ is even exactly when $m$ and $n$ have \textbf{the same parity} (both even or both odd).

\step{2}{Split the $20$ numbers into two groups}
Even group: $2,4,6,\dots,20$ -- $10$ numbers.\quad Odd group: $1,3,5,\dots,19$ -- $10$ numbers.

\step{3}{Count the pairs in the even group}
The first number can be chosen in $10$ ways and the second (different from the first) in $9$ ways: $10\times9=90$ \emph{ordered} choices.
But each pair, for example $\{2,4\}$, is counted twice ($2$ then $4$, or $4$ then $2$), while the condition $m<n$ allows only the order
$(2,4)$. So the number of pairs is $\dfrac{10\times9}{2}=45$.
\begin{recallbox}The number of ways to choose $2$ objects (order does not matter) from $k$ objects is $\dfrac{k\times(k-1)}{2}$.\end{recallbox}

\step{4}{Count the pairs in the odd group and add}
In the same way the odd group gives $45$ pairs. A pair cannot be in both groups, so by the addition principle:
$45+45=90$ pairs.
\begin{warnbox}Forgetting to divide by $2$ gives $90+90=180$ (not an option); forgetting the odd group gives $45$. Always ask yourself:
``Did I count a pair twice? Did I miss a case?''\end{warnbox}
\verify{Count using the complement: the total number of pairs with $m<n$ is $\dfrac{20\times19}{2}=190$. The number of pairs with different
parity (one even, one odd) is $10\times10=100$. So the number of pairs with the same parity is $190-100=90$. It matches!}
So the answer is \textbf{(B)}.
\ans{(B) $90$ pairs}
}

\begin{probbox}{23}{\stars{1}}
Two different two-digit positive integers are called a \emph{reversal pair} if one is obtained from the other by swapping its two digits.
For example, $52$ and $25$ form a reversal pair. The number $2015$ has a special property: it is the product of three different prime
numbers, two of which form a reversal pair ($2015=5\times13\times31$). Including $2015$, how many positive integers less than $10\,000$
have this same property?
\choices{$18$}{$14$}{$20$}{$17$}{$19$}
\end{probbox}
\probsol{
\textbf{Understanding the problem.} We must count the numbers $N<10\,000$ of the form $N=p\times q\times r$, where $p,q,r$ are three
different primes and two of them (call them $p$ and $q$) form a two-digit reversal pair. Plan: find all reversal pairs made of two primes,
then for each pair find the suitable third primes.

\step{1}{Find the reversal pairs in which both numbers are prime}
\begin{recallbox}A prime number is a natural number greater than $1$ whose only divisors are $1$ and itself. The two-digit primes are:
$11, 13, 17, 19, 23, 29, 31, 37, 41, 43, 47, 53, 59, 61, 67, 71, 73, 79, 83, 89, 97$.\end{recallbox}
\begin{tipbox}The units digit of the reversal of $p$ is the tens digit of $p$. A two-digit prime must end in $1,3,7$ or $9$, so we only need
to look at primes starting with $1,3,7,9$: $11,13,17,19,31,37,71,73,79,97$.\end{tipbox}
Check the reversal of each one:
\begin{center}\begin{tabular}{|c|c|c|c|c|c|c|c|c|c|c|}\hline
$p$ & $11$ & $13$ & $17$ & $19$ & $31$ & $37$ & $71$ & $73$ & $79$ & $97$\\\hline
reversal & $11$ & $31$ & $71$ & $91$ & $13$ & $73$ & $17$ & $37$ & $97$ & $79$\\\hline
prime? & same & yes & yes & no & yes & yes & yes & yes & yes & yes\\\hline
\end{tabular}\end{center}
($91=7\times13$ is not prime; $11$ reversed is still $11$, so it does not give two \emph{different} numbers.)
So there are exactly $4$ pairs: $\{13,31\}$, $\{17,71\}$, $\{37,73\}$, $\{79,97\}$.

\step{2}{Compute the product of each pair}
$13\times31=403$;\quad $17\times71=1207$;\quad $37\times73=2701$;\quad $79\times97=7663$.

\step{3}{Find the third prime $r$ for each pair}
We need $r$ prime, different from the two numbers of the pair, and the product $<10\,000$.
\begin{itemize}[leftmargin=1.2em,itemsep=1pt,topsep=2pt]
\item Pair $\{13,31\}$: we need $403\times r<10\,000$. Since $403\times24=9672$ and $403\times25=10\,075$, we need $r\le23$.
The primes $\le23$ are $2,3,5,7,11,13,17,19,23$; remove $13$ (already in the pair): $8$ primes remain, giving
$806;\ 1209;\ 2015;\ 2821;\ 4433;\ 6851;\ 7657;\ 9269$.
\item Pair $\{17,71\}$: we need $1207\times r<10\,000$; since $1207\times8=9656$ and $1207\times9=10\,863$, $r\le8$:
$r\in\{2,3,5,7\}$, giving $2414;\ 3621;\ 6035;\ 8449$ ($4$ numbers).
\item Pair $\{37,73\}$: $2701\times3=8103<10\,000$ and $2701\times4>10\,000$, so $r\in\{2,3\}$: $5402;\ 8103$ ($2$ numbers).
\item Pair $\{79,97\}$: even $7663\times2=15\,326>10\,000$, so there are none.
\end{itemize}

\step{4}{Check that no number is counted twice, and add}
\begin{proofbox}Every number has only one prime factorisation. If a number appeared in two lists, then among its three prime factors there
would be two different reversal pairs; these two pairs would have to share a prime. But the table in step 1 shows that each prime belongs
to exactly one pair. So no number is counted twice.\end{proofbox}
In total: $8+4+2+0=14$ numbers.
\begin{warnbox}Two common mistakes: (1) forgetting that the third prime must be \emph{different} from the two primes of the pair, and so
also counting $403\times13=5239=13\times13\times31$ (wrong, since $13$ is repeated); (2) forgetting the prime $2$.\end{warnbox}
\verify{$2015=403\times5$ is in the list. Checking a few numbers: $9269=13\times31\times23$, $8449=17\times71\times7$,
$8103=37\times73\times3$; all are less than $10\,000$.}
So the answer is \textbf{(B)}.
\ans{(B) $14$}
}

\begin{probbox}{24}{\relax}
The figure is made of $8$ identical small parallelograms. Including these $8$ small parallelograms, how many parallelograms are there in
the figure altogether?
\begin{figC}
  \foreach \j in {0,1,2}{\draw[Ink,thick] (\j*0.55,\j*0.8)--(4.4+\j*0.55,\j*0.8);}
  \foreach \i in {0,...,4}{\draw[Ink,thick] (\i*1.1,0)--(\i*1.1+1.1,1.6);}
\end{figC}
\choices{$29$}{$30$}{$26$}{$27$}{$28$}
\end{probbox}
\probsol{
\textbf{Understanding the problem.} The figure has $2$ rows, each with $4$ small parallelograms. Every parallelogram in the figure is
formed from the segments already drawn: two of its sides lie on \emph{horizontal lines} and two lie on \emph{slanted lines}.

\step{1}{Count the lines in the figure}
There are $3$ horizontal lines (bottom, middle, top) and $5$ slanted lines (parallel to each other).
\begin{figC}
  \fill[tmGreen!30] (1.1,0)--(3.3,0)--(3.85,0.8)--(1.65,0.8)--cycle;
  \foreach \j/\l in {0/1,1/2,2/3}{\draw[tmBlue,thick] (\j*0.55,\j*0.8)--(4.4+\j*0.55,\j*0.8);
     \node[left,text=tmBlue,font=\scriptsize] at (\j*0.55,\j*0.8) {horiz.\ \l};}
  \foreach \i in {0,...,4}{\draw[tmOrange,thick] (\i*1.1,0)--(\i*1.1+1.1,1.6);
     \pgfmathtruncatemacro\k{\i+1}\node[below,text=tmOrange!80!black,font=\scriptsize] at (\i*1.1,0) {s\k};}
\end{figC}
For example, the green parallelogram is determined by horizontal lines $1$ and $2$ and slanted lines s$2$ and s$4$.

\step{2}{Each parallelogram corresponds to a choice of $2$ horizontal lines and $2$ slanted lines}
Choosing $2$ of the $3$ horizontal lines as bottom and top sides, and $2$ of the $5$ slanted lines as left and right sides, gives exactly
one parallelogram; conversely, every parallelogram in the figure comes from exactly one such choice.
\begin{recallbox}The number of ways to choose $2$ lines out of $k$ lines is $\dfrac{k\times(k-1)}{2}$.\end{recallbox}
Choosing $2$ horizontal lines: $\dfrac{3\times2}{2}=3$ ways. \quad Choosing $2$ slanted lines: $\dfrac{5\times4}{2}=10$ ways.

\step{3}{Apply the multiplication principle}
Number of parallelograms: $3\times10=30$.

\step{4}{Check by counting by size}
Call a parallelogram ``height $a$, width $b$'' if it is made of $a$ rows and $b$ columns of small ones:
\begin{center}\begin{tabular}{|c|c|c|c|c|c|}\hline
 & width $1$ & width $2$ & width $3$ & width $4$ & Total\\\hline
height $1$ ($2$ positions in rows) & $2\times4=8$ & $2\times3=6$ & $2\times2=4$ & $2\times1=2$ & $20$\\\hline
height $2$ ($1$ position in rows) & $4$ & $3$ & $2$ & $1$ & $10$\\\hline
\end{tabular}\end{center}
\begin{warnbox}Do not just count the $8$ small ones plus a few ``obvious'' big ones. Count systematically: by size (table above) or by
choosing lines.\end{warnbox}
\verify{The table gives $20+10=30$, the same as $3\times10$.}
So the answer is \textbf{(B)}.
\ans{(B) $30$ parallelograms}
}

\begin{probbox}{25}{\stars{1}}
A jar contains $2$ red marbles and $2$ blue marbles and no other marbles. Angie randomly takes $2$ marbles out of the jar at the same time.
If the two marbles have the same colour, she throws one away and puts the other back into the jar. If they have different colours, she
throws away the red marble and puts the blue marble back into the jar. Angie does this three times in total. The probability that the last
marble left in the jar is red is
\choices{$\dfrac12$}{$\dfrac14$}{$\dfrac23$}{$\dfrac13$}{$0$}
\end{probbox}
\probsol{
\textbf{Understanding the problem.} Each draw removes exactly $1$ marble from the jar. After $3$ draws, $1$ of the $4$ marbles is left.
We follow the \emph{state} of the jar, written (number of red, number of blue), starting from $(2\text{R},2\text{B})$.

\step{1}{First draw, from $(2\text{R},2\text{B})$}
Call the marbles R$_1$, R$_2$, B$_1$, B$_2$. There are $6$ equally likely ways to take $2$ marbles:
R$_1$R$_2$;\ B$_1$B$_2$;\ R$_1$B$_1$;\ R$_1$B$_2$;\ R$_2$B$_1$;\ R$_2$B$_2$.
\begin{itemize}[leftmargin=1.2em,itemsep=1pt,topsep=2pt]
\item Two reds ($1$ way): throw away $1$ red $\to(1\text{R},2\text{B})$.
\item Two blues ($1$ way): throw away $1$ blue $\to(2\text{R},1\text{B})$.
\item Different colours ($4$ ways): throw away the red $\to(1\text{R},2\text{B})$.
\end{itemize}
So after draw $1$: state $(1\text{R},2\text{B})$ with probability $\frac56$, state $(2\text{R},1\text{B})$ with probability $\frac16$.

\step{2}{Second draw}
From $(1\text{R},2\text{B})$: $3$ ways to draw; $2$ ways give different colours $\to(0\text{R},2\text{B})$; $1$ way gives two blues $\to(1\text{R},1\text{B})$.\\
From $(2\text{R},1\text{B})$: $3$ ways to draw; $1$ way gives two reds $\to(1\text{R},1\text{B})$; $2$ ways give different colours $\to(1\text{R},1\text{B})$.
So from $(2\text{R},1\text{B})$ we certainly go to $(1\text{R},1\text{B})$.

\step{3}{Third draw}
From $(0\text{R},2\text{B})$: two blues are drawn, one is thrown away $\to(0\text{R},1\text{B})$.\\
From $(1\text{R},1\text{B})$: different colours are drawn, the red is thrown away $\to(0\text{R},1\text{B})$.
\begin{figC}
  \tikzset{st/.style={draw=tmBlue,fill=tmBlue!8,rounded corners=3pt,inner sep=3pt,font=\footnotesize}}
  \node[st] (a) at (0,0) {$2$R, $2$B};
  \node[st] (b) at (2.8,0.9) {$1$R, $2$B};\node[st] (c) at (2.8,-0.9) {$2$R, $1$B};
  \node[st] (d) at (5.6,0.9) {$0$R, $2$B};\node[st] (e) at (5.6,-0.9) {$1$R, $1$B};
  \node[st,draw=AnsRed,fill=AnsRed!8] (f) at (8.2,0) {$0$R, $1$B};
  \draw[->,thick,Slate] (a)--node[above,font=\scriptsize]{$\frac56$}(b);
  \draw[->,thick,Slate] (a)--node[below,font=\scriptsize]{$\frac16$}(c);
  \draw[->,thick,Slate] (b)--node[above,font=\scriptsize]{$\frac23$}(d);
  \draw[->,thick,Slate] (b)--node[above right,font=\scriptsize,pos=0.3]{$\frac13$}(e);
  \draw[->,thick,Slate] (c)--node[below,font=\scriptsize]{$1$}(e);
  \draw[->,thick,Slate] (d)--node[above,font=\scriptsize]{$1$}(f);
  \draw[->,thick,Slate] (e)--node[below,font=\scriptsize]{$1$}(f);
  \node[font=\scriptsize,text=Slate] at (1.4,-1.7) {draw 1};\node[font=\scriptsize,text=Slate] at (4.2,-1.7) {draw 2};
  \node[font=\scriptsize,text=Slate] at (6.9,-1.7) {draw 3};
\end{figC}
Every branch ends at $(0\text{R},1\text{B})$: the last marble is \textbf{always} blue.

\step{4}{Why -- without drawing the tree}
\begin{proofbox}Look at the number of blue marbles. Whenever at least one of the two marbles drawn is blue, Angie always puts a blue marble back
into the jar (two blues: she puts one back; different colours: she puts the blue back). The number of blue marbles decreases only when two
blues are drawn, and then the one that was put back is still there. So the jar \emph{always contains at least one blue marble}. When only one
marble is left, it must be blue.\end{proofbox}
Therefore the event ``the last marble is red'' is an \emph{impossible event}, with probability $0$.
\begin{recallbox}An impossible event (one that never happens) has probability $0$; a certain event has probability $1$.\end{recallbox}
\begin{warnbox}Seeing ``$2$ red, $2$ blue'' and immediately guessing $\frac12$ is a trap: the throwing-away rule is \emph{not symmetric}
in the two colours (with different colours, the red one is always thrown away).\end{warnbox}
\verify{Total probability of reaching the final state $(0\text{R},1\text{B})$: $\frac56\cdot\frac23+\frac56\cdot\frac13+\frac16\cdot1=\frac{10}{18}+\frac{5}{18}+\frac{3}{18}=1$.}
So the answer is \textbf{(E)}.
\ans{(E) $0$}
}

\needspace{22\baselineskip}
\begin{probbox}{26}{\stars{1}}
On the $10\times10$ grid of squares shown, triangle $FGH$ is drawn. Point $P$ can be any of the $41$ points where grid lines cross
that lie \emph{inside} (not on the sides of) triangle $FGH$. For each position of $P$, exactly three triangles are formed: $FPG$, $GPH$,
$HPF$. Among all $41\times3=123$ of these triangles, how many have an area equal to exactly half the area of triangle $FGH$?
\begin{figC}\begin{scope}[scale=0.3]
  \draw[Line,thin] (0,0) grid (10,10);
  \draw[Ink,line width=1.3pt] (0,0)--(10,0)--(5,10)--cycle;
  \node[above,font=\scriptsize] at (5,10) {$F$};\node[below left,font=\scriptsize] at (0,0) {$G$};
  \node[below right,font=\scriptsize] at (10,0) {$H$};
\end{scope}\end{figC}
\choices{$5$}{$7$}{$3$}{$11$}{$9$}
\end{probbox}
\probsol{
\textbf{Understanding the problem.} Take the side of one small square as the unit and put $G$ at the bottom-left corner. Then $G$ and $H$
are the two bottom corners ($GH=10$) and $F$ is the midpoint of the top side. Each grid point is named by two numbers (squares to the
right, squares up); for example $G=(0,0)$, $H=(10,0)$, $F=(5,10)$.

\step{1}{Area of triangle $FGH$ and half of it}
\begin{recallbox}Area of a triangle $=\dfrac12\times\text{base}\times\text{corresponding height}$.\end{recallbox}
$\text{Area}(FGH)=\frac12\times10\times10=50$. Half of it is $25$.

\step{2}{Triangle $GPH$: the easiest}
Triangle $GPH$ has base $GH=10$, and its height is the distance from $P$ down to $GH$, that is the ``height'' $h$ of $P$ (the number of
squares above the bottom). $\text{Area}(GPH)=\frac12\times10\times h=5h$. We need $5h=25$, so $h=5$: $P$ lies on the $5$th horizontal grid
line. At height $5$ (half the height) the triangle is half as wide as its base, from $2.5$ to $7.5$. The grid points strictly \emph{inside}
on that segment are $(3,5),(4,5),(5,5),(6,5),(7,5)$: $\mathbf{5}$ points.

\step{3}{Triangle $FPG$: using a parallel line}
\begin{recallbox}Triangles with a common base whose third vertices lie on a line parallel to the base have the same height, so they have the
same area.\end{recallbox}
Let $M=(5,0)$ be the midpoint of $GH$. Triangle $FMG$ has base $GM=5$ and height $10$, so $\text{Area}(FMG)=\frac12\times5\times10=25$.
So every point $P$ on the line through $M$ parallel to $FG$ gives $\text{Area}(FPG)=25$.
\begin{proofbox}If $P$ is farther from $FG$ than that line, the height of triangle $FPG$ (on base $FG$) is larger, so the area is more than
$25$; if $P$ is closer to $FG$, the area is less than $25$. So $\text{Area}(FPG)=25$ \emph{only when} $P$ lies on that line.\end{proofbox}
Side $FG$ goes from $(0,0)$ to $(5,10)$: for every $1$ square to the right it goes up $2$ squares. The parallel line through $M=(5,0)$ does the
same: it passes through $(5,0)$, $(6,2)$, $(7,4)$ and then meets side $FH$ at $(7.5,\,5)$ (the midpoint of $FH$). The grid points
\emph{inside} the triangle are $(6,2)$ and $(7,4)$: $\mathbf{2}$ points. ($M$ lies on side $GH$, so it is excluded.)

\step{4}{Triangle $HPF$: by symmetry}
The figure is symmetric about the vertical line through $F$, so similarly there are $\mathbf{2}$ points: $(4,2)$ and $(3,4)$.
\begin{figC}\begin{scope}[scale=0.42]
  \draw[Line,thin] (0,0) grid (10,10);
  \draw[Ink,line width=1.2pt] (0,0)--(10,0)--(5,10)--cycle;
  \draw[tmBlue,line width=1.2pt] (2.5,5)--(7.5,5);
  \draw[tmOrange,line width=1.2pt] (5,0)--(7.5,5);\draw[tmGreen!80!black,line width=1.2pt] (5,0)--(2.5,5);
  \foreach \p in {(3,5),(4,5),(5,5),(6,5),(7,5)}{\fill[tmBlue] \p circle (5pt);}
  \foreach \p in {(6,2),(7,4)}{\fill[tmOrange] \p circle (5pt);}
  \foreach \p in {(4,2),(3,4)}{\fill[tmGreen!80!black] \p circle (5pt);}
  \node[above,font=\scriptsize] at (5,10) {$F$};\node[below left,font=\scriptsize] at (0,0) {$G$};
  \node[below right,font=\scriptsize] at (10,0) {$H$};\node[below,font=\scriptsize] at (5,0) {$M$};
  \node[right,font=\scriptsize,text=tmBlue,align=left] at (10.4,6) {blue: area $GPH=25$};
  \node[right,font=\scriptsize,text=tmOrange!80!black] at (10.4,4.6) {orange: area $FPG=25$};
  \node[right,font=\scriptsize,text=tmGreen!70!black] at (10.4,3.2) {green: area $HPF=25$};
\end{scope}\end{figC}

\step{5}{Add}
Number of triangles with area $25$: $5+2+2=9$. (One point $P$ cannot give two triangles of area $25$, because then the third triangle would
have area $50-25-25=0$, while $P$ lies strictly inside. So no triangle is missed or counted twice.)
\verify{For $P=(6,2)$: triangle $FPG$ fits in the rectangle from $(0,0)$ to $(6,10)$, of area $60$. Cut off three right triangles:
$\frac12\times6\times2=6$, $\frac12\times1\times8=4$ and $\frac12\times5\times10=25$. What remains is $60-6-4-25=25$. Correct!}
So the answer is \textbf{(E)}.
\ans{(E) $9$ triangles}
}

\begin{probbox}{27}{\relax}
A $10\times10$ square grid is made of $100$ points, as shown. Point $P$ is marked. One of the other $99$ points is chosen at random to be
$Q$. The probability that segment $PQ$ is horizontal or vertical is
\begin{figC}\begin{scope}[scale=0.38]
  \foreach \x in {0,...,9}{\foreach \y in {0,...,9}{\fill[Navy] (\x,\y) circle (3.2pt);}}
  \fill[white] (6,3) circle (6pt);\draw[tmOrange,line width=1.2pt] (6,3) circle (6pt);
  \node[above right,font=\scriptsize,text=tmOrange!80!black] at (6,3) {$P$};
\end{scope}\end{figC}
\choices{$\dfrac2{11}$}{$\dfrac15$}{$\dfrac1{10}$}{$\dfrac4{25}$}{$\dfrac5{33}$}
\end{probbox}
\probsol{
\textbf{Understanding the problem.} $Q$ is chosen at random among $99$ points (not counting $P$), all equally likely.
Segment $PQ$ is horizontal when $Q$ is in the same row as $P$, and vertical when $Q$ is in the same column as $P$.

\step{1}{Number of possible outcomes}
There are $100-1=99$ ways to choose $Q$.

\step{2}{Number of favourable outcomes}
The row containing $P$ has $10$ points; without $P$ there are $9$. The column containing $P$ has $10$ points; without $P$ there are $9$.
A point other than $P$ cannot be in both the same row and the same column as $P$ (the only such point is $P$ itself), so no point is
counted twice: $9+9=18$ points in total.
\begin{figC}\begin{scope}[scale=0.3]
  \foreach \x in {0,...,9}{\foreach \y in {0,...,9}{\fill[Slate!50] (\x,\y) circle (3pt);}}
  \foreach \x in {0,...,9}{\fill[tmBlue] (\x,3) circle (5pt);}
  \foreach \y in {0,...,9}{\fill[tmGreen!80!black] (6,\y) circle (5pt);}
  \fill[white] (6,3) circle (6pt);\draw[tmOrange,line width=1.2pt] (6,3) circle (6pt);
  \node[right,font=\scriptsize,text=tmBlue] at (10,3) {$9$ points in the same row};
  \node[above,font=\scriptsize,text=tmGreen!70!black] at (6,9.3) {$9$ points in the same column};
\end{scope}\end{figC}

\step{3}{Compute the probability}
\[P=\frac{18}{99}=\frac{18\div9}{99\div9}=\frac{2}{11}.\]
\begin{warnbox}The denominator is $99$, not $100$, because $Q$ must be different from $P$. Using $100$ gives $\frac{18}{100}$, which is not
an option. Also do not count $10+10=20$ points (that counts $P$ twice).\end{warnbox}
\begin{tipbox}The position of $P$ in the grid does not matter: every point has exactly $9$ points in its row and $9$ points in its column.\end{tipbox}
\verify{$\frac2{11}\times99=18$, exactly the number of favourable points.}
So the answer is \textbf{(A)}.
\ans{(A) $\dfrac2{11}$}
}

\begin{probbox}{28}{\stars{1}}
The sum of all the digits of the integers from $98$ to $101$ is
\[9+8+9+9+1+0+0+1+0+1=38.\]
What is the sum of all the digits of the integers from $1$ to $2008$?
\choices{$30\,054$}{$27\,018$}{$28\,036$}{$30\,036$}{$28\,054$}
\end{probbox}
\probsol{
\textbf{Understanding the problem.} Writing out $2008$ numbers and adding is far too long. We split them into three blocks: $1$ to $999$;
$1000$ to $1999$; $2000$ to $2008$, and find the digit sum of each block by counting \emph{how many times each digit appears}.

\step{1}{Block $1$ to $999$: write leading zeros to get three digits}
Write the numbers from $0$ to $999$ with three digits: $000, 001, 002, \dots, 998, 999$. Adding leading zeros and the number $000$
\emph{does not change} the digit sum. There are $1000$ numbers, that is $3000$ digits.
\begin{proofbox}In the units place the digits repeat in the cycle $0,1,2,\dots,9$; $1000$ numbers contain $100$ cycles, so each digit
$0$--$9$ appears exactly $100$ times. In the tens place each digit appears $10$ times in a row ($00$--$09$ have tens digit $0$, \dots)
and this repeats $10$ times: again $100$ times. In the hundreds place: $000$--$099$ have hundreds digit $0$ ($100$ times), \dots,
$900$--$999$ have hundreds digit $9$ ($100$ times). So in each place, each digit appears $100$ times.\end{proofbox}
Hence each digit $0,1,\dots,9$ appears $3\times100=300$ times. The digit sum is
\[300\times(0+1+2+\dots+9)=300\times45=13\,500.\]
\begin{recallbox}$0+1+2+\dots+9=\dfrac{9\times10}{2}=45$ (pair them up: $0+9=1+8=\dots=4+5=9$, and there are $5$ pairs).\end{recallbox}

\step{2}{Block $1000$ to $1999$}
Each number has the form $1\,abc$ where $abc$ runs from $000$ to $999$. The last three digits give $13\,500$ (as in step 1). The thousands
digit is $1$, appearing $1000$ times, which adds $1000$. Total: $13\,500+1000=14\,500$.

\step{3}{Block $2000$ to $2008$}
There are $9$ numbers. The thousands digit $2$ contributes $9\times2=18$; the two middle digits are $0$; the units digits are
$0,1,\dots,8$, contributing $0+1+\dots+8=36$. Total: $18+36=54$.

\step{4}{Add the three blocks}
\[13\,500+14\,500+54=28\,054.\]
\begin{warnbox}Do not forget the last small block $2000$--$2008$ (forgetting it gives $28\,000$), and do not forget the thousands digit $1$
in the second block (forgetting it gives $27\,054$).\end{warnbox}
\verify{Test the method on a small case, $1$ to $99$: each digit appears $2\times10=20$ times, sum $20\times45=900$. Directly: the units
place contributes $10\times45=450$ and the tens place contributes $10\times(1+2+\dots+9)=450$, total $900$. It matches.}
So the answer is \textbf{(E)}.
\ans{(E) $28\,054$}
}

\begin{probbox}{29}{\relax}
Two fair dice are rolled, one red and one blue (each has six faces numbered $1$ to $6$). The probability that the number on the red die
is greater than the number on the blue die is
\choices{$\dfrac{18}{36}$}{$\dfrac{25}{36}$}{$\dfrac{15}{36}$}{$\dfrac{12}{36}$}{$\dfrac{17}{36}$}
\end{probbox}
\probsol{
\textbf{Understanding the problem.} An outcome is a pair (number on the red die, number on the blue die). The dice have different
colours, so $(5,2)$ and $(2,5)$ are different outcomes.

\step{1}{Number of possible outcomes}
By the multiplication principle: $6\times6=36$ equally likely outcomes.

\step{2}{Make a table and mark the cells ``red $>$ blue''}
\begin{figC}\begin{scope}[scale=0.5]
  \foreach \r in {1,...,6}{\foreach \b in {1,...,6}{
    \pgfmathsetmacro\yy{6-\r}
    \ifnum\r>\b \fill[AnsRed!25] (\b-1,\yy) rectangle ++(1,1);
    \else\ifnum\r=\b \fill[Slate!20] (\b-1,\yy) rectangle ++(1,1);\fi\fi}}
  \draw[Ink] (0,0) grid (6,6);
  \foreach \i in {1,...,6}{\node[font=\scriptsize,text=tmBlue] at (\i-0.5,6.4) {\i};
     \node[font=\scriptsize,text=AnsRed] at (-0.4,6.5-\i) {\i};}
  \node[font=\scriptsize,text=tmBlue] at (3,7.1) {blue die};
  \node[font=\scriptsize,text=AnsRed,rotate=90] at (-1.2,3) {red die};
\end{scope}\end{figC}
Red cells: the red number is greater than the blue number; grey cells: the two numbers are equal. Count row by row (red number):
red $=1$: $0$ cells; red $=2$: $1$ cell; red $=3$: $2$ cells; red $=4$: $3$ cells; red $=5$: $4$ cells; red $=6$: $5$ cells.
Total: $0+1+2+3+4+5=15$ favourable outcomes.

\step{3}{Compute the probability}
$P(\text{red}>\text{blue})=\dfrac{15}{36}$.

\step{4}{A quick way: use symmetry}
There are $6$ outcomes where the two numbers are equal (the diagonal). The other $36-6=30$ outcomes split equally between ``red is
greater'' and ``blue is greater'' (swap the roles of the dice), so each kind has $30\div2=15$ outcomes.
\begin{warnbox}The answer $\frac{18}{36}=\frac12$ is a trap: it forgets the cases where the two numbers are equal. ``Red is greater'' and
``blue is greater'' are not the only possibilities.\end{warnbox}
\verify{$P(\text{red}>\text{blue})+P(\text{equal})+P(\text{blue}>\text{red})=\frac{15}{36}+\frac6{36}+\frac{15}{36}=1$.}
So the answer is \textbf{(C)}.
\ans{(C) $\dfrac{15}{36}$}
}

\begin{probbox}{30}{\relax}
Two fair dice are rolled, one red and one blue. The probability that the \emph{product} of the two numbers shown is even is
\choices{$\dfrac12$}{$\dfrac34$}{$\dfrac14$}{$\dfrac23$}{$\dfrac{11}{36}$}
\end{probbox}
\probsol{
\textbf{Understanding the problem.} As in problem 29, there are $6\times6=36$ equally likely outcomes. Counting the even products directly
is rather long, so we count the \emph{complement}: the outcomes with an odd product.

\step{1}{When is the product of two numbers odd?}
\begin{recallbox}odd $\times$ odd $=$ odd;\quad even $\times$ (any number) $=$ even. So the product of two numbers is odd exactly when
\emph{both} numbers are odd.\end{recallbox}

\step{2}{Count the outcomes with an odd product}
Odd numbers on each die: $1,3,5$ ($3$ numbers). Both dice odd: $3\times3=9$ outcomes.

\step{3}{Count the outcomes with an even product using the complement}
\begin{recallbox}\emph{Complement}: number of favourable outcomes $=$ number of possible outcomes $-$ number of unfavourable outcomes.
Likewise, $P(\text{event})=1-P(\text{event does not happen})$.\end{recallbox}
Number of outcomes with an even product: $36-9=27$.

\step{4}{Compute the probability}
\[P(\text{even product})=\frac{27}{36}=\frac{27\div9}{36\div9}=\frac34.\]
\begin{warnbox}Many students think ``even and odd are half and half'' and choose $\frac12$. But just \emph{one} even number makes the product
even, so an even product happens much more often.\end{warnbox}
\verify{Direct count: red die even ($3$ ways) with any blue number ($6$ ways): $18$ outcomes; red die odd ($3$ ways) and blue die even
($3$ ways): $9$ outcomes. Total $18+9=27$. It matches.}
So the answer is \textbf{(B)}.
\ans{(B) $\dfrac34$}
}
